\documentclass[10pt,a4paper]{amsart}
\usepackage[margin=19mm]{geometry}
\usepackage{amsmath,amssymb,mathtools}
\usepackage[scr=rsfs,cal=euler]{mathalfa}
\usepackage{xfrac}
\usepackage[unicode,hidelinks,bookmarks=false]{hyperref}
\newcommand{\ie}{i.e.\ }
\newcommand{\cf}{cf.\ }
\newcommand{\resp}{respectively\ }
\newcommand{\Subsection}{\subsection}
\providecommand{\nobreakdash}{\nobreak\hbox{-}}
\setlength{\emergencystretch}{2em}
\multlinegap0pt
\usepackage{algorithm}\usepackage{algpseudocode}
\usepackage{tikz}
\newtheorem{theorem}{Theorem}
\newcommand{\KL}{D_{\mathrm{KL}}}
\newcommand{\dd}{\mathop{}\mathopen{}\mathrm{d}}
\title{\texorpdfstring{Exposing the Illusion of Fairness: \\ Auditing Vulnerabilities to Distributional Manipulation Attacks}{Exposing the Illusion of Fairness: - Auditing Vulnerabilities to Distributional Manipulation Attacks}}
\author{Valentin Lafargue, Adriana Laurindo Monteiro, Emmanuelle Claeys, Laurent Risser, Jean-Michel Loubes}
\date{Source: arXiv:2507.20708v3 (CC BY 4.0)}
\begin{document}
\maketitle

\noindent\emph{Source and licence.} \texorpdfstring{Exposing the Illusion of Fairness: \\ Auditing Vulnerabilities to Distributional Manipulation Attacks}{Exposing the Illusion of Fairness: - Auditing Vulnerabilities to Distributional Manipulation Attacks}. Version arXiv:2507.20708v3 (CC BY 4.0), distributed by arXiv under the Creative Commons Attribution 4.0 International licence, http://creativecommons.org/licenses/by/4.0/, as recorded in the arXiv licence metadata for this exact version and shown on the abstract page https://arxiv.org/abs/2507.20708v3. Full licence notice: \texttt{licenses/arxiv-2507.20708-CC-BY-4.0.txt}.\par\smallskip

\noindent\textit{Editorial scope.} Counted unit: the article's theorem thm:discrete:reweigth:multidim (source0.tex lines 331--357). The article's complete original proof of it is delivered with it (source0.tex lines 1180--1188, one \texttt{proof} environment of 9 lines, which the article places away from the statement). No mathematics is written, repaired or re-derived: the statement and the proof are the article's own lines, in the article's own notation, and no retained line is substituted or re-flowed.  No further source-local context is needed: every cross-reference of every retained line resolves inside the two delivered spans.  Nothing was omitted from any retained span, so the package carries no omission record.  No retained line contains a citation, so the wrapper carries no bibliography and no bibliography entry.  No retained line contains a figure environment, an image-inclusion command or drawing code, so no image is delivered and the package declares no asset dependency.  Editorial formatting only, in addition: an AMS \texttt{amsart} preamble that restates the article's own declarations and macros used by the retained lines, namely the environments \texttt{theorem} on separate counters, together with the standard packages those lines need.  The retained environments print in this document as Theorem 1 in order of appearance, whereas the article prints the same items with the numbering of its own full document.  The complete native source of the article as distributed in the arXiv e-print for this exact version is bundled unchanged as \texttt{sources/182307/source0.tex}.
\begin{theorem}[Entropic Projection under constraint] \label{thm:discrete:reweigth:multidim}
	Let $t \in \mathbb{R}^k$ and $\Phi : E \to \mathbb{R}^k$ be measurable.
	Assume that $t$ can be written as a convex combination of $\Phi(X_1,S_1, \widehat{Y}_1,Y_1) , \ldots , \Phi(X_n,S_n, \widehat{Y}_n,Y_n)$, with positive weights. Assume also that the empirical covariance matrix of $\Phi(X,S,\hat{Y}, Y)$ is invertible.
	% $\mathbb{E}_{Q_n} (\Phi \Phi^\top)  - \mathbb{E}_{Q_n} (\Phi ) \mathbb{E}_{Q_n} (\Phi^\top)$ is invertible.

	Let $\mathcal{D}_{\Phi,t}$ be the set of all probability measures $P$ on $E$ such that $
	\int_{E} \Phi(x) \dd P(x)=t$.
		Then, \begin{equation} \label{thesol:Qn}
Q_t:=\mathrm{arginf}_{P\in\mathcal{D}_{\Phi,t}}\KL(P \Vert Q_n)
\end{equation}
    exists and is unique. It can also be computed as 

        \begin{equation}
            Q_t=
		\frac{1}{n}
		\sum_{i=1}^n
		\lambda_i^{(t)}
		\delta_{X_i , S_i, \widehat{Y}_i , Y_i}
        \end{equation}
        
        with, for $i\in\{1,\ldots,n\}$,
	$\lambda^{(t)}_i
		= \exp
		\bigg( \langle\xi(t) ,  \Phi(X_i,S_i, \widehat{Y_i},Y_i) \rangle - \log Z(\xi(t))
		\bigg)$, where for a  vector $\xi \in \mathbb{R}^k$, let $Z(\xi):=\frac{1}{n} \sum_{i=1}^n e^{\langle  \Phi(X_i,S_i, \widehat{Y}_i,Y_i) , \xi \rangle}$; and $\xi(t)$ is the unique minimizer of the strictly convex function
	$H(\xi):=\log Z(\xi)-\langle\xi,t\rangle$.
\end{theorem}

\begin{proof}
    
Theorem~\ref{thm:discrete:reweigth:multidim} implies the existence of a distribution $Q_t$ such that \[DI(f,Q_t,S)= \frac{\lambda_0 + \delta_0}{\lambda_1 - \delta_1}\frac{n_1}{n_0}=t_1. \]
 We have
\[\Delta_{DI} =
 \frac{n_1}{n_0}\bigg( \frac{\lambda_1 \delta_0 + \lambda_0 \delta_1)}{(\lambda_1 - \delta_1)\lambda_1}\bigg)
\] 
Among all possible solutions, we privilege the two solutions described in the Proposition. Knowing the new DI desired, we can obtain a set of solution for $\delta_0$ and $\delta_1$.
 \end{proof}

\end{document}
